\section*{Chapter \ref{ch:g12:saving-and-investment} --- Saving, Investment and the Financial System}
% figdata: figdata/grade-12/07-saving-and-investment/out/ (si, si_key, crisis1929, crisis2007, crisis_key); Brenland invented.

\begin{solution}{exo:g12:saving-and-investment:1}
$S = 1000 - 600 - 200 = 200$, and $I = S = 200$.
\end{solution}

\begin{solution}{exo:g12:saving-and-investment:2}
$S = 1000 - 650 - 200 = 150$; \omterm{def:g10:national-accounts:net-exports}{net exports} $S - I = -30$.
\end{solution}

\begin{solution}{exo:g12:saving-and-investment:3}
Singapore saves about 46\,\% and invests 27\,\%: it lends the difference abroad, with a surplus of \omterm{def:g10:national-accounts:net-exports}{net exports}. Mozambique invests about 42\,\% and saves 12\,\%: it borrows the
difference or receives it as investment from abroad, with a deficit of \omterm{def:g10:national-accounts:net-exports}{net exports}.
\end{solution}

\begin{solution}{exo:g12:saving-and-investment:4}
Private: $1000 - 250 - 600 = 150$. Public: $250 - 280 = -30$. National: $150 - 30 = 120$.
\end{solution}

\begin{solution}{exo:g12:saving-and-investment:5}
Households $+80$, firms $-60$; together $+20$, which these two sectors lend to the others.
\end{solution}

\begin{solution}{exo:g12:saving-and-investment:6}
$100 + 20r = 180 - 20r$ gives $r = 2\,\%$, and $S = I = 140$.
\end{solution}

\begin{solution}{exo:g12:saving-and-investment:7}
$100 + 20r = 200 - 20r$: $r = 2.5\,\%$, $S = I = 150$.
\end{solution}

\begin{solution}{exo:g12:saving-and-investment:8}
\omterm{def:g12:saving-and-investment:national-saving}{National saving} $80 + 20r = 180 - 20r$: $r = 2.5\,\%$; investment 130; \omterm{def:g12:saving-and-investment:national-saving}{private saving} 150. Investment falls by 10: half the deficit is crowded out, and higher \omterm{def:g12:saving-and-investment:national-saving}{private saving}
pays for the other half.
\end{solution}

\begin{solution}{exo:g12:saving-and-investment:9}
$140 + 20r = 180 - 20r$: $r = 1\,\%$, investment 160. In an open economy at a fixed world rate, investment would not change; the extra saving would be lent abroad, raising \omterm{def:g10:national-accounts:net-exports}{net
exports} by 40.
\end{solution}

\begin{solution}{exo:g12:saving-and-investment:10}
$350/500 = 70\,\%$.
\end{solution}

\begin{solution}{exo:g12:saving-and-investment:11}
1933: 69.2 and 24.9\,\%. 2009: 94.5 and 9.3\,\%. The first was far deeper: output per person fell by almost a third, against about 5\,\%.
\end{solution}

\begin{solution}{exo:g12:saving-and-investment:12}
In a closed economy each country's investment equals its saving, so all points would lie on the line of equal rates. Across 155 countries they scatter widely around it, and
investment rises by only about 0.25 points per point of saving: saving and investment are separated by borrowing and lending across borders.
\end{solution}

\begin{solution}{exo:g12:saving-and-investment:13}
\omterm{def:g12:present-value:present-value}{Present value} $5/0.05 = 100$. A buyer at 150 receives 5 and expects a gain of 15: $20/150 \approx 13.3\,\%$, more than 5\,\%, so buying looks right while the expectation lasts.
When buyers stop expecting the rise, the price has no support above 100 and can fall by a third or more.
\end{solution}

\begin{solution}{exo:g12:saving-and-investment:14}
It receives $0.8 \times 30 = 24$ for assets worth 30: a loss of 6, leaving capital of 2. Other banks holding the same assets must value them at the lower price, and lose too;
some may then have to sell, pushing the price down further.
\end{solution}

\begin{solution}{exo:g12:saving-and-investment:15}
At a fixed world rate, investment depends only on that rate. The government borrows from abroad: foreign lenders buy its \omterm{def:g12:present-value:bond}{bonds}, and Brenland runs a deficit of \omterm{def:g10:national-accounts:net-exports}{net
exports}. In return it owes interest, and eventually repayment, to foreigners.
\end{solution}

\begin{solution}{pb:g12:saving-and-investment:1}
\textbf{1.} $16 + 2r = 24 - 2r$: $r = 2\,\%$.

\textbf{2.} 20\,\% of GDP each: $0.2 \times 393\,264 \approx 78\,653$ million dalens.

\textbf{3.} \omterm{def:g12:saving-and-investment:national-saving}{Public saving} 0; \omterm{def:g12:saving-and-investment:national-saving}{national saving} 20\,\% of GDP.

\textbf{4.} A project is worth undertaking when the \omterm{def:g12:present-value:present-value}{present value} of its returns exceeds its cost; a lower rate raises that \omterm{def:g12:present-value:present-value}{present value}, so more projects pass.

\textbf{5.} Each dalen saved buys more future consumption.

\textbf{6.} $S = 16 + 2r - 2 = 14 + 2r$.

\textbf{7.} $14 + 2r = 24 - 2r$: $r = 2.5\,\%$.

\textbf{8.} Investment $24 - 5 = 19\,\%$; \omterm{def:g12:saving-and-investment:national-saving}{private saving} $16 + 5 = 21\,\%$.

\textbf{9.} By 1 point of GDP: a crowding-out ratio of $1/2 = 0.5$.

\textbf{10.} Households and firms, whose saving rises by 1 point at the higher rate.

\textbf{11.} Investment $24 - 4 = 20\,\%$; \omterm{def:g12:saving-and-investment:national-saving}{national saving} $16 + 4 - 2 = 18\,\%$.

\textbf{12.} $18 - 20 = -2\,\%$ of GDP: \omterm{def:g10:national-accounts:net-exports}{net exports} of $-0.02 \times 393\,264 \approx -7865$ million dalens.

\textbf{13.} Lenders in Kestria and elsewhere, who buy Brenland's \omterm{def:g12:present-value:bond}{bonds} or assets.

\textbf{14.} Closed: investment falls by 1 point. Open: it does not fall; the deficit is matched by borrowing from abroad.

\textbf{15.} Interest at the world rate on the borrowing, about $0.02 \times 7865 \approx 157$ million dalens a year in real terms for each year of such borrowing, unless it is
repaid.

\textbf{16.} The open economy's: saving and investment can differ because capital flows, though the loose link shows flows are not perfectly free.

\textbf{17.} $24 - 6 = 18\,\%$ of GDP.

\textbf{18.} Its borrowing is a large part of the world's demand for funds, so it moves the world rate as the closed economy's deficit moved its own.

\textbf{19.} If the borrowing pays for public investment that raises future output, and the taxes on that output, by more than the interest owed.

\textbf{20.} In the closed economy half the deficit crowds out investment: a ratio of 0.5. Once capital flows at a world rate of 2\,\%, Brenland runs a deficit of \omterm{def:g10:national-accounts:net-exports}{net exports}
of 2\,\% of GDP, about 7865 million dalens.
\end{solution}
